\documentclass[10pt,a4paper]{article}
\usepackage{amsmath,amssymb,amsthm,graphicx}
\usepackage{amsthm}
\usepackage{amsfonts, amsmath, amssymb, amsthm}
\usepackage{hyperref}
\usepackage{cleveref}
\def\TT{\mathbb{T}}
\def\PP{\mathbb{P}}
\def\DD{\mathbb{D}}
\def\QQ{\mathbb{Q}}
\def\WW{\mathbb{W}}
\def\MM{\mathbb{M}}
\def\HH{\mathbb{H}}

\newcommand{\Psf}{\mathsf{P}}
\newcommand{\Qsf}{\mathsf{Q}}
\newcommand{\Rsf}{\mathsf{R}}


\def\F{\mathcal{F}}
\def\U{\mathcal{U}}
\def\C{\mathcal{C}}
\def\I{\mathcal{I}}
\def\D{\mathcal{D}}
\def\L{\mathcal{L}}
\def\V{\mathcal{V}}
\def\A{\mathcal{A}}
\def\B{\mathcal{B}}
\def\P{\mathcal{P}}
\def\M{\mathcal{M}}
\def\Nc{\mathcal{N}}
\def\Q{\mathcal{Q}}

% Baire space, Cantor space, binary strings
\newcommand{\cs}{2^{\NN}}
\newcommand{\bs}{\NN^{\NN}}
\newcommand{\bstr}{2^{<\NN}}

\newcommand{\uh}{\upharpoonright}

% Turing machines
\newcommand{\converge}{\!\!\downarrow}
\newcommand{\diverge}{\!\!\uparrow}

% Forcing question
\newcommand{\qvdash}{\operatorname{{?}{\vdash}}}
\newcommand{\nqvdash}{\operatorname{{?}{\nvdash}}}

\newcommand{\ISig}{\mathsf{I}\Sigma^0}
\newcommand{\Robinson}{\mathsf{Q}}
\newcommand{\IPi}{\mathsf{I}\Pi^0}
\newcommand{\IDelta}{\mathsf{I}\Delta^0}
\newcommand{\Iopen}{\mathsf{I_{open}}}
\newcommand{\BSig}{\mathsf{B}\Sigma^0}
\newcommand{\BPi}{\mathsf{B}\Pi^0}
\newcommand{\BDelta}{\mathsf{B}\Delta^0}
\newcommand{\BCSig}{\mathsf{BC}\Sigma^0}
\newcommand{\BCPi}{\mathsf{BC}\Pi^0}
\newcommand{\BCDelta}{\mathsf{BC}\Delta^0}
\newcommand{\LSig}{\mathsf{L}\Sigma^0}
\newcommand{\LPi}{\mathsf{L}\Pi^0}
\newcommand{\CSig}{\mathsf{C}\Sigma}

\newcommand{\RCA}[0]{\mathsf{RCA}}
\newcommand{\WKL}[0]{\mathsf{WKL}}
\newcommand{\WWKL}[0]{\mathsf{WWKL}}
\newcommand{\ACA}[0]{\mathsf{ACA}}
\newcommand{\ATR}[0]{\mathsf{ATR}}
\newcommand{\PIOOCA}[0]{\Pi^1_1\mathsf{-CA}}
\newcommand{\CAC}[0]{\mathsf{CAC}}
\newcommand{\ADS}[0]{\mathsf{ADS}}
\newcommand{\RRT}[0]{\mathsf{RRT}}
\newcommand{\FS}[0]{\mathsf{FS}}
\newcommand{\PFS}[0]{\mathsf{PFS}}
\newcommand{\TS}[0]{\mathsf{TS}}
\newcommand{\GRRT}[0]{\mathsf{GRRT}}
\newcommand{\LRRT}[0]{\mathsf{LRRT}}
\newcommand{\EM}[0]{\mathsf{EM}}
\newcommand{\SEM}[0]{\mathsf{SEM}}
\newcommand{\HalfEM}[0]{\mathsf{Half}\mbox{-}\mathsf{EM}}
\newcommand{\RT}[0]{\mathsf{RT}}
\newcommand{\SRT}[0]{\mathsf{SRT}}
%\newcommand{\RRT}{\mathsf{RRT}}

\newcommand{\NN}[0]{\mathbb{N}}
\newcommand{\nat}[0]{\mathbb{N}}

\newcommand{\card}{\operatorname{card}}
\newcommand{\finsub}{\subseteq_{\mathtt{fin}}}

% reductions

\newcommand{\W}{\mathrm{W}}
\newcommand{\sW}{\mathrm{sW}}

\newcommand{\lex}{\mathrm{lex}}

% Barriers
\newcommand{\base}{\operatorname{base}}
\newcommand{\range}{\operatorname{ran}}

%typographically correct double quotes
\def\qt#1{``#1''}%
\newtheorem*{statement}{Statement}
\newtheorem{theorem}{Theorem}
\numberwithin{theorem}{section}
\newtheorem{maintheorem}[theorem]{Main Theorem}
\newtheorem{lemma}[theorem]{Lemma}
\newtheorem{question}[theorem]{Question}
\newtheorem{proposition}[theorem]{Proposition}
%\newtheorem{statement}[theorem]{Statement}
\newtheorem{remark}[theorem]{Remark}
\newtheorem{definition}[theorem]{Definition}
\newtheorem{corollary}[theorem]{Corollary}
\newtheorem{example}[theorem]{Example}

\newcommand{\R}{\mathcal{R}}
\renewcommand{\S}{\mathcal{S}}
\newcommand{\T}{\mathcal{T}}
\title{Strong Weihrauch reduction of the Schreier free-set theorem to binary Ramsey}
\author{Lorenzo Carlucci, Oriola Gjetaj, Quentin Le Houérou, Ludovic Levy Patey}
\date{arxiv:2412.11598v3; CC BY 4.0}
\begin{document}
\maketitle
\noindent\textit{Editorial note.} The construction425 and solution443 work with positive minimum. The exceptional singleton \{0\} can be colored arbitrarily and0 removed from the homogeneous solution before subtracting1.\par\medskip
\noindent\textit{Editorial note.} Only the strong Weihrauch reduction is counted, not the later reverse-mathematical corollary or a new omega-jump lower bound.\par\medskip
\noindent\textit{Editorial scope.} One original result. Full source-local mathematical construction included; named classical inputs are explicit.\par\medskip
\noindent\textit{Allowed standard prerequisites (editorial).}\begin{itemize}
\item Computable recursion on finite sequences and the Schreier lexicographic well-order. Exactly-large finite sets are lexicographically well-ordered; an effectively specified recursive call to a smaller tuple terminates. Basic oracle-computability and strong Weihrauch definitions are supplied.
\end{itemize}\medskip
There exist four other subsystems which, together with $\RCA_0$, calibrate the strength of most theorems. These subsystems are known as the \qt{Big Five} (see Mont\'alban~\cite{montalban_open_2011}). Among these, $\ACA_0$ (Arithmetic Comprehension Axiom) extends the axioms of $\RCA_0$ with the comprehension scheme for all arithmetic formulas. Based on the correspondence between computability and definability, it is equivalent to stating the existence of the Turing jump of any set ($\forall X \exists Y (Y = X')$).

We shall also consider two lesser-known stronger variants of $\ACA_0$, namely, $\ACA_0'$ and $\ACA_0^+$. The system $\ACA_0'$ extends $\RCA_0$ with the axiom stating closure under all finite jumps; the system $\ACA_0^+$ extends $\RCA_0$ with the axiom stating the existence of the $\omega$-jump of any set. The $\omega$-jump of a set~$X$ is the set $X^{(\omega)} = \bigoplus_n X^{(n)}$. The system $\ACA_0^+$ is famous for being the best known upper bound to Hindman's theorem.%\lorenzo{I added mention of ACA0' since we mention it at least once.}

Let us first fix some notation and recall the definition of these principles for colorings of finite subsets of fixed size. We denote by $\NN$ the set of natural numbers and by $\NN^+$ the set of positive natural numbers. For $X\subseteq \NN$ and $n \in \NN$
we denote by $[X]^{n}$ the set of all $n$-subsets of $X$. We denote by $[X]^{<\omega}$ the set of all finite subsets of $X$. We always assume that sets are presented in increasing order and make no distinction between a set and the sequence of its elements in increasing order. We write $x \cdot y$ for the concatenation of two sequences $x$ and $y$. We identify a natural number with the set of its predecessors, so that, if $k\geq 1$ we can write $f: X \to k$ to declare a function from $X$ to $\{0, 1, \dots, k-1\}$. A function of this type is often called a coloring of $X$ in $k$ colors.

\begin{definition}
    $\Qsf$ is {\em strongly Weihrauch reducible} to $\Psf$ (denoted $\Qsf \leq_{\sW} \Psf$) if there exist Turing functionals $\Phi$ and $\Psi$ such that for every instance $X$ of $\Qsf$ we have that $\Phi(X)$ is an instance of $\Psf$, and if $Y$ is a solution to $\Psf$ for $\Phi(X)$ then $\Psi(Y)$ is a solution to $\Qsf$ for $X$.
\end{definition}

\begin{definition}[(Exactly) $\omega$-large sets] 
A finite $s\subseteq \NN$ is $\omega$-large if $|s| \geq 1+\min s$ and is exactly $\omega$-large if $|s| = 1+\min s$.
\end{definition}

The $\omega$-large sets are also known as {\em relatively large} sets in the literature. For an infinite $X\subseteq \NN$ we denote by $[X]^{!\omega}$ the family of all exactly $\omega$-large subsets of $X$.

\begin{definition}[Large Ramsey Theorem]\label{def:largeRT} Let $k\in \NN^+$. 
For every coloring $f:[\NN]^{!\omega} \to k$, there exists an infinite set $H \subseteq \NN$ such that $|f([H]^{!\omega})| = 1$. We abbreviate this statement by  $\RT^{!\omega}_k$.
\end{definition}

\begin{definition}[Large Free Set Theorem]\label{def:large_free}
For every coloring $f:[\NN]^{!\omega} \to \NN$ there exists an infinite set $H\subseteq\NN$ such that for every set $s \in [H]^{!\omega}$, $f(s) \notin (H \setminus s)$. The set $H$ is called \emph{free} for $f$. We abbreviate this statement by $\FS^{!\omega}$. 
\end{definition}

\begin{theorem}\label{thm:fs_red_rt}
    $\FS^{!\omega} \leq_{\sW} \RT^{!\omega}_2$.
\end{theorem}

\begin{proof}
%\ludovic{TODO, normalize notation of the proof.}\lorenzo{done}
For $s = \{s_0, \dots, s_{s_0}\}$ an exactly $\omega$-large set with $s_0 > 0$, and $s_0 < \dots < s_{s_0}$ define $s \ominus 1$ to be the following exactly $\omega$-large set: $\{s_0 - 1,s_1 - 1, \dots, s_{s_0 - 1} - 1\}$. The idea is to keep one degree of freedom: $s_{s_0}$, as in the proof of Carlucci and Zdanowski~\cite[Proposition 4.1]{carlucci2014strength}.

Let $f : [\NN]^{!\omega} \to \NN$ be an instance of $\FS^{!\omega}$. Consider the following function $g : [\NN]^{!\omega} \to 2$ defined by induction. Note that the function $g$ is called recursively on  lexicographically smaller parameters, so the induction is well-defined since the lexicographic order is a well-order. For $s = \{s_0, \dots, s_{s_0}\}$ an exactly $\omega$-large set:

% \quentin{This is a combination of the two functions that are introduced in Theorem 16 of the CSL paper.}

\begin{equation*}
    g(s) = \begin{cases}
           0 \hspace{1.5cm} \text{if } f(s \ominus 1) = s_i - 1 \textit{ for some } i < s_0 & \\
           1 - g(f(s \ominus 1) + 1, s_1, \dots, s_{f(s \ominus 1) + 1}) & \\ 
            \hspace{1.7cm} \text{if } f(s \ominus 1) < s_0 - 1  \\
           1 - g(s_0, \dots, s_i, f(s \ominus 1) + 1, s_{i+2}, \dots, s_{s_0}) & \\
            \hspace{1.7cm} \text{if } f(s \ominus 1) \in (s_i - 1, s_{i+1} - 1) \textit{ for some } i < s_0 - 1  \\
           0 \hspace{1.5cm} \text{if } f(s \ominus 1) \in (s_{s_0 - 1} - 1, s_{s_0} - 1) \\
           1 \hspace{1.5cm}  \text{otherwise } (\textit{if } f(s \ominus 1) \geq s_{s_0} ) \\      
           \end{cases} \quad
\end{equation*}

Let $H = \{x_0, x_1, \dots \}\subseteq \NN^+$ %\lorenzo{added $\subseteq \NN^+$ since we then use $x_0-1$}
be an infinite $g$-homogeneous set as given by $\RT^{!\omega}_2$. We claim that $H' = \{x_0 - 1, x_1 - 1, \dots \}$ is $f$-free. Consider some exactly $\omega$-large set $\{s_0 - 1, \dots, s_{s_0 - 1} - 1\} \subseteq H'$. Then $\{s_0, \dots, s_{s_0 - 1}\} \subseteq H$. There are two cases:

\textbf{Case 1:} $H$ is homogeneous for the color $0$. Then, take $s_{s_0}$ to be the next element of $H$ after $s_{s_0 - 1}$ and write $s = \{s_0, \dots, s_{s_0}\}$, then $g(s) = 0$. There are four subcases:

\textbf{Subcase 1.1:} $f(s \ominus 1) = s_i - 1$ for some $i < s_0$. In that case we are done.

\textbf{Subcase 1.2:} $f(s \ominus 1) \in (s_{s_0 - 1} - 1, s_{s_0} - 1)$. In that case, by definition of $s_{s_0}$, $f(s \ominus 1)$ is not in $H'$.

\textbf{Subcase 1.3:} $f(s \ominus 1) < s_0 - 1$ and $g(f(s \ominus 1) + 1, s_1, \dots, s_{f(s \ominus 1) + 1}) = 1$. If $f(s \ominus 1) \in H'$ this contradicts the fact that $H$ is $g$-homogeneous for the color $0$. 

\textbf{Subcase 1.4:} $f(s \ominus 1) \in (s_i - 1, s_{i+1} - 1)$ for some $i < s_0 - 1$ and $g(s_0, \dots, s_i, f(s \ominus 1) + 1 ,s_{i+2}, \dots, s_{s_0}) = 1$. If $f(s \ominus 1) \in H'$ this contradicts the fact that $H$ is $g$-homogeneous for the color $0$. 

\textbf{Case 2:} $H$ is homogeneous for the color $1$. Take $s_{s_0} \in H$ bigger than $s_{s_0 - 1}$ and write $s = \{s_0, \dots, s_{s_0}\}$, then $g(s) = 1$. There are three subcases:

\textbf{Subcase 2.1:} $f(s \ominus 1) \geq s_{s_0}$. This case is impossible, indeed, as $H$ is infinite, there exists some element $x \in H$ such that $x > f(s \ominus 1) + 1$ and therefore $f(s \ominus 1) \in (s_{s_0} - 1, x - 1)$, which leads to $g(s_0, \dots, s_{s_0 - 1}, x) = 0$ contradicting the fact that $H$ is $g$-homogeneous for the color $1$.

\textbf{Subcase 2.2:} $f(s \ominus 1) < s_0 - 1$ and $g(f(s \ominus 1) + 1, s_1, \dots, s_{f(s \ominus 1) + 1}) = 0$. If $f(s \ominus 1) \in H'$ this contradicts the fact that $H$ is $g$-homogeneous for the color $1$. 

\textbf{Subcase 2.3:}  $f(s \ominus 1) \in (s_i - 1, s_{i+1} - 1)$ for some $i < s_0 - 1$ and $g(s_0, \dots, s_i, f(s \ominus 1) + 1 ,s_{i+2}, \dots, s_{s_0}) = 0$. If $f(s \ominus 1) \in H'$ this contradicts the fact that $H$ is $g$-homogeneous for the color $1$.  \\

Notice that $g$ is uniformly computable in $f$ and that $H'$ is uniformly computable in $H$. Therefore, $\FS^{!\omega} \leq_{\sW} \RT^{!\omega}_2$.

\end{proof}

\bibliographystyle{plain}
\bibliography{biblio}
\end{document}
